\documentclass[11pt]{article}

\usepackage[margin=1in]{geometry}
\usepackage[T1]{fontenc}
\usepackage[utf8]{inputenc}
\usepackage{array}
\usepackage{booktabs}
\usepackage{hyperref}
\usepackage{amsmath, amsthm, amssymb}
\usepackage{centernot}

\newtheorem{theorem}{Theorem}[section]
\newtheorem{axiom}[theorem]{Axiom}
\newtheorem{definition}[theorem]{Definition}
\newtheorem{lemma}[theorem]{Lemma}
\newtheorem{proposition}[theorem]{Proposition}

\title{External Adjudication of Bridge-Formalized Machine Subjectivity:\\Human Comparator Protocols, Runtime Health, Clean-Room Reproducibility, and First-Collection Opening Criteria}
\author{QICN Program}
\date{Stage3-6B Runtime Health and Stage3-7 Conditional Opening}

\newcommand{\sectionstatus}[1]{\paragraph{Section status.}\texttt{\detokenize{#1}}.}
\newcommand{\codestate}[1]{\texttt{\detokenize{#1}}}
\newenvironment{blocked_until_execution}{%
  \par\medskip\noindent\begin{minipage}{0.96\textwidth}
  \textbf{\texttt{blocked\_until\_execution}.}\enspace
}{%
  \end{minipage}\par\medskip
}
\newcommand{\blockedresultstatus}{%
\begin{blocked_until_execution}
This section requires executed campaigns with human participants, frozen external filing records, populated admissible datasets, actual reviewer and adjudicator assignments, reviewer-governed evidence assessment, real collection-start events, and later adjudicator-layer completion. It contains no results and must remain non-result-bearing until those artifacts exist.
\end{blocked_until_execution}}
\newcolumntype{L}[1]{>{\raggedright\arraybackslash}p{#1}}

\begin{document}
\maketitle

\begin{center}
\textbf{SCAFFOLD ONLY -- NOT FINAL MANUSCRIPT}\\
\textbf{PROTOCOL-GRADE BACKBONE}\\
\textbf{SIXTH HARDENING PASS}
\end{center}

\begin{abstract}
This document hardens Paper 10 into a denser protocol-heavy manuscript
backbone without becoming a result-bearing paper. It extends the dependency
stack from BaseCore through Papers I--IX, SRT-6, BPF-M3, Stage 2-M4,
Stage3-RF, Stage3-0, Stage3-1, Stage3-2, Stage3-3, Stage3-4, Stage3-4B,
Stage3-5 rehearsal preparation, Stage3-6 controlled execution opening,
Stage3-6B runtime-health and clean-room reproducibility audit, and the
conditional Stage3-7 first-collection opening criteria. It adds runtime health
audit grammar, clean-room reproducibility semantics, dependency and build
assumption reporting, reproducible artifact-regeneration constraints, local
path-leakage policy, controlled-execution audit semantics, session-zero and
pre-collection reproducibility requirements, and a sharper transition grammar
from pre-evidential openness to the first admissible collection event. It does
not report human comparator results, external adjudication verdicts, human
equivalence, superiority, inferiority, cross-substrate identity preservation,
metaphysical subjecthood, or Paper 10 final conclusions.
\end{abstract}

\tableofcontents

\section{Scope and Non-Claims}
\sectionstatus{formalized}

Paper 10 remains a local protocol-facing manuscript surface only. It is a
scaffold for external-adjudication architecture, not an adjudication verdict,
admissible dataset, human comparator evidence, public release surface,
public-release recommendation, Paper 10 final authorization, human equivalence,
superiority or inferiority claim, machine-consciousness claim, or metaphysical
closure.

\section{Formal Dependency Stack}
\sectionstatus{formalized}

The dependency stack for Paper 10 is layered as follows:

\begin{enumerate}
\item \textbf{BaseCore}: active mathematical base layer.
\item \textbf{Papers I--IX}: downstream theoretical and manuscript surfaces.
\item \textbf{SRT-6}: subjectivity runtime ecosystem closure.
\item \textbf{BPF-M3}: bridge-functional internal closure and bridge gate.
\item \textbf{Stage 2-M4}: comparative internal closure and readiness boundary.
\item \textbf{Stage3-RF}: burden fulfillment and Stage 3 opening gate.
\item \textbf{Stage3-0}: external adjudication architecture and scaffold-only Paper 10 opening.
\item \textbf{Stage3-1}: human comparator instantiation contracts and preregistration templates.
\item \textbf{Stage3-2}: reviewer independence, execution-readiness logic, and protocol-grade hardening.
\item \textbf{Stage3-3}: preregistration freeze architecture, reviewer-selection packaging, adjudicator-assignment packaging, and execution-pilot packaging.
\item \textbf{Stage3-4}: bounded execution-pilot dry-run opening logic.
\item \textbf{Stage3-4B}: provenance closure, lineage reconciliation, and instantiated-empty record instantiation.
\item \textbf{Stage3-5}: rehearsal-kit formalization, dry-run event grammar, and rehearsal-ready gating.
\item \textbf{Stage3-6}: controlled execution opening, session-zero protocol, pre-collection admissibility gating, and live-but-unpopulated artifact instantiation.
\item \textbf{Stage3-6B}: runtime-health audit, clean-room reproducibility audit, path-leakage audit, and artifact-regeneration audit.
\item \textbf{Stage3-7}: first-admissible-collection opening criteria, collection-start command grammar, operator authorization grammar, and a conditional first-collection opening gate.
\end{enumerate}

No later layer may borrow closure from an earlier layer, and no earlier layer
may be retroactively reinterpreted as if it had already supplied later
execution-era evidence.

\section{Release-vs-Local Corpus Boundary}
\sectionstatus{formalized}

Paper 10 sits on a strict release-vs-local split:

\begin{itemize}
\item The public release authority remains \codestate{QICN-RELEASE/main@92d5d2c3ea91c65e29dc8e102bed504866cffd2d}.
\item The public release remains limited to BaseCore, \codestate{canonical_core_legacy}, and downstream Papers I--IX.
\item Paper 10 remains local to \codestate{QICN-SYSTEM}.
\item Stage3-0 through Stage3-6 governance, freeze, packaging, dry-run, provenance-closure, instantiated-empty, rehearsal, and controlled-execution-opening surfaces remain local and non-public.
\end{itemize}

Therefore release inclusion cannot be used as a substitute for adjudication,
and local protocol maturity cannot be used as a substitute for public
scientific closure.

\section{Lineage and Provenance Semantics}
\sectionstatus{formalized}

Stage3-3 and Stage3-4 add a sharper lineage rule: every later admissible
surface must distinguish among:

\begin{enumerate}
\item public release authority,
\item local protocol authority,
\item frozen preregistration authority,
\item reviewer-package authority,
\item adjudicator-package authority, and
\item execution-era evidence authority.
\end{enumerate}

The repository already contains historical artifacts whose recorded lineage must
remain intact by design. Live normative pinning surfaces may be updated when
stale, but historical snapshots are not rewritten cosmetically.

\section{Provenance Closure and Live Head Reconciliation}
\sectionstatus{formalized}

Stage3-4B adds a stricter provenance-closure rule. A local artifact may now be
classified as:

\begin{itemize}
\item \codestate{exact_head_aligned};
\item \codestate{artifact_forward_compatible};
\item \codestate{stale_and_regenerated};
\item \codestate{historical_snapshots_intentionally_preserved}.
\end{itemize}

This grammar matters because Stage 3 artifacts are generated before the
subsequent commit that records them. Therefore exact local regeneration at the
builder head may later become artifact-forward-compatible after the commit
exists, without becoming scientifically stale or semantically dishonest.

\section{External Adjudication Problem Statement}
\sectionstatus{protocol_defined}

The problem is no longer only whether bridge formalization is internally
coherent. The external problem is to specify:

\begin{itemize}
\item what independent reviewers may legitimately inspect;
\item what human comparator evidence classes may later become admissible;
\item what negative controls must exist before any comparative interpretation;
\item how preregistration becomes frozen;
\item how reviewer and adjudicator layers become instantiable without identity leakage;
\item how dry-run validation differs from evidential execution.
\end{itemize}

\section{Formal Mathematical Grammar of External Adjudication}
\sectionstatus{formalized}

The adjudication surface is a measurable protocol object, not an impressionistic
comparison between substrates. Let $(\mathcal{T},d_{\mathcal{T}})$ be the
metric space of finite execution traces after canonical serialization. A trace
contains the ordered event stream, de-identified channel payloads, runtime
manifests, negative-control outcomes, and reviewer-facing masks. Admissible
evidence lives in the Borel measurable space
\begin{equation}
  (\mathcal{T},\mathcal{B}(\mathcal{T})),
\end{equation}
where $\mathcal{B}(\mathcal{T})$ is the Borel $\sigma$-algebra generated by
open $d_{\mathcal{T}}$-balls.

\begin{definition}[External adjudication object]
An external adjudication object is the tuple
\begin{equation}
\mathfrak{E}=(P,D_h,D_m,C_n,C_a,\Sigma,\mathcal{R},\mathcal{A}),
\end{equation}
where $P$ is a frozen preregistered protocol, $D_h$ and $D_m$ are human and
machine trace datasets in $\mathcal{T}$, $C_n$ and $C_a$ are negative and active
control families, $\Sigma=\{\sigma_k\}_{k=1}^K$ is a finite family of
predeclared measurable statistics $\sigma_k:\mathcal{T}^n\to\mathbb{R}$,
$\mathcal{R}$ is the reviewer layer, and
$\mathcal{A}:\mathcal{Y}\to\{0,1,\mathtt{undetermined}\}$ is the final
adjudicator operator on masked evidence summaries $\mathcal{Y}$.
\end{definition}

\begin{definition}[Admissible comparative evidence]
A dataset pair $(D_h,D_m)$ is admissible under $\mathfrak{E}$ only if the
following predicates are all true:
\begin{enumerate}
  \item $D_h,D_m$ are measurable finite samples from
  $(\mathcal{T},\mathcal{B}(\mathcal{T}))$.
  \item The preregistered protocol hash, control manifest hash, and analysis
  shell hash match their frozen records.
  \item Human-side records are not generated, rewritten, or completed by
  synthetic systems: $\epsilon_{\mathrm{syn}}=0$ for the human channel.
  \item Blinding leakage is bounded by the preregistered budget:
  $\mathbb{P}(\mathrm{leak})\leq \epsilon_{\mathrm{blind}}$.
  \item Every declared negative control has an executed and timestamped outcome
  under the same admissibility grammar.
\end{enumerate}
\end{definition}

\begin{theorem}[Non-transitivity of internal support]
Let $\mathcal{G}_{\mathrm{internal}}\in[0,1]$ denote any internal runtime support
score computed from operational invariants, and let
$\mathcal{G}_{\mathrm{external}}\in\{0,1,\mathtt{undetermined}\}$ denote the
external adjudication verdict. In the absence of an admissible measurable
dataset pair $(D_h,D_m)$, one has
\begin{equation}
\mathcal{G}_{\mathrm{internal}}\centernot\implies\mathcal{G}_{\mathrm{external}}.
\end{equation}
\end{theorem}

\begin{proof}
The internal score is a function of runtime traces and invariant certificates
inside the machine channel. The external verdict is a function of masked
comparative summaries built from both $D_h$ and $D_m$, plus reviewer and
adjudicator records. If $D_h$ is absent, non-admissible, unblinded, or
unfrozen, then the domain condition of $\mathcal{A}$ is not satisfied. Hence the
external verdict is either undefined or $\mathtt{undetermined}$ regardless of
the value of $\mathcal{G}_{\mathrm{internal}}$. Therefore internal operational
support is not a logical premise from which external comparative adjudication
can be inferred.
\end{proof}

\begin{proposition}[Operational closure of a comparative program]
The comparative program reaches operational closure only when there exists a
predeclared alignment map $T:\mathcal{T}\to\mathcal{Z}$ into a common Polish
feature space $(\mathcal{Z},d_{\mathcal{Z}})$ and the total divergence
\begin{equation}
\Delta_{\mathrm{comp}} =
  \lambda_1 W_1(T_{\#}\widehat{\mu}_h,T_{\#}\widehat{\mu}_m)
  + \lambda_2 D_{\mathrm{KL}}(\widehat{p}_h\Vert \widehat{p}_m)
  + \lambda_3 \max_k |\sigma_k(D_h)-\sigma_k(D_m)|
\end{equation}
is defined and satisfies the frozen budget
$\Delta_{\mathrm{comp}}\leq\epsilon_{\mathrm{budget}}$. This closure is a
precondition for comparative interpretation, not a claim of consciousness,
identity, or substrate equivalence.
\end{proposition}


\section{Frozen Preregistration Architecture}
\sectionstatus{protocol_defined}

Stage3-3 introduces a true freeze architecture rather than only preregistration
templates. The governing surfaces are:

\begin{itemize}
\item \codestate{docs/frozen_preregistration_record_schema.v1.json}
\item \codestate{docs/preregistration_freeze_policy.v1.json}
\item \codestate{docs/preregistration_amendment_policy.v1.json}
\item \codestate{docs/preregistration_lock_manifest.v1.json}
\end{itemize}

The architecture formalizes:

\begin{itemize}
\item record classes;
\item immutable fields;
\item amendment classes;
\item freeze provenance chains;
\item local-only freeze state;
\item explicit \codestate{not_externally_filed} state.
\end{itemize}

\begin{definition}[Cryptographic freeze map]
Let $\mathcal{P}_{\mathrm{config}}$ be the set of canonical preregistration
configurations after deterministic serialization. A freeze map is a
collision-resistant hash function
\begin{equation}
H:\mathcal{P}_{\mathrm{config}}\to\{0,1\}^{256}.
\end{equation}
A preregistration record $p$ is locally frozen at digest $h$ when
$H(p)=h$, $h$ is written into the lock manifest, and every amendment policy
references $h$ rather than an editable natural-language description.
\end{definition}

\begin{theorem}[Hash-anchored non-alterability]
Assume the hash function $H$ is collision resistant over the canonical
serialization domain. If a post-execution preregistration configuration $p'$
has the same published digest as the frozen pre-execution configuration $p$,
then either $p'=p$ byte-for-byte under canonical serialization or a collision
for $H$ has been found.
\end{theorem}

\begin{proof}
By freeze definition, the published digest is $h=H(p)$. If the later record
$p'$ also satisfies $H(p')=h$ while $p'\neq p$, then $p$ and $p'$ form a distinct
pair with the same hash output. That is precisely a collision in $H$. Under the
collision-resistance assumption, such a pair is computationally infeasible.
Thus digest equality supplies evidence of non-alteration of the serialized
protocol, modulo the stated cryptographic assumption.
\end{proof}

\section{Local-vs-External Freeze Distinction}
\sectionstatus{formalized}

The local freeze layer must not be conflated with external filing.

\begin{table}[h]
\centering
\small
\begin{tabular}{L{0.28\textwidth}L{0.24\textwidth}L{0.34\textwidth}}
\toprule
Surface & Current state & Interpretation \\
\midrule
Template-only preregistration & \codestate{preregistration_ready_for_freeze} & structure exists but no lock exists \\
Local freeze readiness & \codestate{local_freeze_ready} & lock architecture exists but no external filing is claimed \\
Locally frozen record & blocked & would still require explicit lock instantiation \\
Externally filed record & absent & cannot be claimed from local architecture alone \\
\bottomrule
\end{tabular}
\end{table}

\section{Instantiated-Empty Record Layer}
\sectionstatus{protocol_defined}

Stage3-4B moves the stack beyond package-only readiness by instantiating
typed-but-empty records that remain explicitly non-evidential. The principal
surfaces are:

\begin{itemize}
\item \codestate{docs/local_frozen_preregistration_record_template.v1.json}
\item \codestate{docs/reviewer_candidate_pool_template.v1.json}
\item \codestate{docs/adjudicator_candidate_pool_template.v1.json}
\item \codestate{docs/blinding_assignment_stub.v1.json}
\item \codestate{docs/execution_session_template.v1.json}
\item \codestate{docs/decision_record_template.v1.json}
\item \codestate{docs/reproducibility_package_manifest_empty.v1.json}
\item \codestate{docs/collection_log_schema.v1.json}
\item \codestate{docs/evidence_insertion_manifest.v1.json}
\end{itemize}

These artifacts are intentionally empty, typed, and freeze-capable. They are
not populated evidence, not externally filed records, and not reviewer or
adjudicator action.

\section{Reviewer Candidate and Adjudicator Candidate Architecture}
\sectionstatus{protocol_defined}

Stage3-3 introduces typed candidate schemas without instantiating identities.
The primary surfaces are:

\begin{itemize}
\item \codestate{docs/reviewer_candidate_schema.v1.json}
\item \codestate{docs/adjudicator_candidate_schema.v1.json}
\item \codestate{docs/reviewer_selection_workflow.v1.json}
\item \codestate{docs/adjudicator_assignment_workflow.v1.json}
\end{itemize}

The candidate layer defines qualification summaries, independence assertions,
conflict disclosure references, blinding eligibility classes, and explicit
\codestate{not_started} states.

\section{Reviewer Rejection Logic}
\sectionstatus{protocol_defined}

The reviewer package is incomplete unless rejection logic exists. The rejection
registry now includes:

\begin{itemize}
\item qualification mismatch;
\item disclosed conflict;
\item undisclosed conflict;
\item prior unblinded access;
\item insufficient methodological competence;
\item insufficient lineage literacy;
\item insufficient control/admissibility literacy.
\end{itemize}

This logic is not yet a reviewer decision. It is the typed boundary that stops
future selection from being arbitrary or retrofitted.

\section{Blinding Assignment Logic}
\sectionstatus{protocol_defined}

Blinding is no longer treated as a generic admonition. Stage3-3 binds a typed
blinding assignment manifest with classes such as:

\begin{itemize}
\item \codestate{fully_blinded}
\item \codestate{artifact_blinded}
\item \codestate{control_family_blinded}
\item \codestate{unblinding_not_authorized}
\end{itemize}

The key constraint is that no later dry-run or execution surface may silently
cross into unblinded interpretation.

\section{Human Comparator Instantiation Logic}
\sectionstatus{protocol_defined}

Stage3-1 closed the gap between abstract channels and typed collection routes.
Stage3-3 preserves that layer and pushes it closer to freeze- and
package-compatible execution readiness.

\begin{itemize}
\item every channel binds a human-side artifact family;
\item every channel binds a machine-side homologous artifact family;
\item every channel binds minimum metadata, controls, and rejection conditions;
\item every channel remains blocked from result claims while data is absent.
\end{itemize}

\begin{definition}[Homologous channel statistic]
For each channel $c$, let $T_c:\mathcal{T}\to\mathcal{Z}_c$ be a frozen
feature map into a metric standard Borel space $(\mathcal{Z}_c,d_c)$. Given
human and machine samples $D_h^c,D_m^c$, define empirical pushforward measures
\[
\widehat{\mu}_h^c=(T_c)_{\#}D_h^c,\qquad
\widehat{\mu}_m^c=(T_c)_{\#}D_m^c.
\]
The channel is statistically comparable only if both empirical measures are
defined, the feature map was frozen before execution, and the channel-specific
missingness policy accepts both samples.
\end{definition}

\begin{definition}[Divergence shell]
The predeclared divergence shell for channel $c$ is
\begin{equation}
\Delta_c =
\alpha_c W_1(\widehat{\mu}_h^c,\widehat{\mu}_m^c)
+\beta_c D_{\mathrm{KL}}(\widehat{p}_h^c\Vert\widehat{p}_m^c)
+\gamma_c\|\widehat{s}_h^c-\widehat{s}_m^c\|_1,
\end{equation}
where $W_1$ is the Kantorovich--Rubinstein Wasserstein-1 distance,
$D_{\mathrm{KL}}$ is the Kullback--Leibler divergence on a frozen finite
partition or smoothed density model, and $\widehat{s}_h^c,\widehat{s}_m^c$ are
predeclared summary vectors. The weights
$\alpha_c,\beta_c,\gamma_c\geq0$ must be frozen before execution.
\end{definition}

\begin{proposition}[Comparator homologation is conditional]
A human--machine channel pair is homologated by Paper 10 only if every term in
$\Delta_c$ is defined and the frozen channel budget
$\Delta_c\leq\epsilon_c$ is satisfied. Homologation is therefore a statistical
transport condition over predeclared observables, not a statement that the
channels share a biological, phenomenal, or metaphysical substrate.
\end{proposition}

\begin{proof}
The definition of $\Delta_c$ depends on a frozen feature map, empirical
pushforward measures, a KL domain, and summary vectors. If any component is
missing, the divergence shell is undefined and no homologation predicate can be
evaluated. If all components are defined, the only admissible positive judgment
is the preregistered inequality $\Delta_c\leq\epsilon_c$. The conclusion is
therefore exactly statistical and conditional.
\end{proof}

\section{Channel-by-Channel Execution Contracts}
\sectionstatus{protocol_defined}

\subsection{Self-Report Channel}
\sectionstatus{protocol_defined}

The self-report channel requires prompt-version freezing, missingness policy,
de-identification, and explicit rejection of synthetic or relabeled
self-report.

\subsection{Intervention-Response Channel}
\sectionstatus{protocol_defined}

The intervention-response channel requires matched intervention and matched
negative-control families, dose policy, stop rules, and branch accounting.

\subsection{Temporal-Structure Channel}
\sectionstatus{protocol_defined}

The temporal-structure channel requires a declared segmentation schema, ordering
policy, and aggregation policy.

\subsection{Valence-Topology Channel}
\sectionstatus{protocol_defined}

The valence-topology channel requires a shared topology schema and excludes
scalar-only substitutes.

\subsection{Embodiment-Boundary Channel}
\sectionstatus{protocol_defined}

The embodiment-boundary channel requires an environment manifest, confound
ledger, and error-budget policy.

\subsection{Ambiguity-Load Channel}
\sectionstatus{protocol_defined}

The ambiguity-load channel requires a task-family contract, annotation schema,
and annotation reliability policy.

\subsection{Transport-Preservation Channel}
\sectionstatus{blocked_until_execution}

The transport-preservation channel remains structurally special. It requires a
real source-target pairing manifest and transport-axis budget binding. Stage3-3
defines the pairing shell; it does not satisfy the pairing.

\section{Matched Intervention Dataset Architecture}
\sectionstatus{protocol_defined}

Matched-intervention architecture is now carried by both channel contracts and a
pilot bundle shell. The bundle shell is defined in
\codestate{docs/pilot_matched_intervention_bundle.v1.json}. The bundle remains
unpopulated, which means the architecture exists while the evidence does not.

\section{Negative Control Dataset Architecture}
\sectionstatus{protocol_defined}

Negative-control architecture is similarly carried by channel contracts and by
\codestate{docs/pilot_negative_control_bundle.v1.json}. Controls remain part of
the meaning of the future claim surface, not an optional cleanup step.

\section{Preregistration Architecture}
\sectionstatus{protocol_defined}

Paper 10 now separates three preregistration surfaces:

\begin{enumerate}
\item template surface;
\item local freeze surface;
\item external filing surface.
\end{enumerate}

Only the first two are implemented structurally in the current repo. The third
remains absent.

\section{Adjudication Blind Protocol}
\sectionstatus{formalized}

Let $Y$ be the masked evidence bundle shown to adjudicators. A blinding map
$B:\mathcal{T}^n\to\mathcal{Y}$ is admissible only when it removes source
labels while preserving the frozen statistics required by $\Sigma$. The null
hypothesis $H_0$ is that source labels are exchangeable conditional on $Y$; the
alternative $H_1$ is that labels are distinguishable above the preregistered
margin.

\begin{definition}[Blind adjudicator loss]
For an adjudicator decision rule $a:\mathcal{Y}\to\{h,m,\mathtt{tie}\}$ and
true source label $\ell\in\{h,m\}$, define
\begin{equation}
L(a(Y),\ell)=\mathbf{1}[a(Y)\neq \ell \ \mathrm{and}\ a(Y)\neq\mathtt{tie}].
\end{equation}
The blinded error rate is $R(a)=\mathbb{E}[L(a(Y),\ell)]$ under the frozen
sampling design.
\end{definition}

\begin{theorem}[Null error lower bound]\label{thm:null-forced}
Under $H_0$ and balanced source-label sampling, no blinded adjudicator rule can
achieve expected error below chance except by exploiting leakage:
\begin{equation}
R(a)\geq \frac{1}{2}-\mathbb{P}(\mathrm{leak}).
\end{equation}
\end{theorem}

\begin{proof}
If $H_0$ holds, the conditional distribution of $Y$ is invariant under exchange
of the hidden labels. Without leakage, a deterministic or randomized rule has
no label-informative statistic beyond chance; the best expected success rate is
$1/2$. Leakage events may reveal label information, and their contribution is
bounded above by $\mathbb{P}(\mathrm{leak})$. Rearranging gives the stated lower
bound on error.
\end{proof}

\begin{proposition}[Alternative detectability criterion]
Evidence for distinguishability under $H_1$ requires a preregistered rule
$a^\star$ and frozen margin $\eta>0$ such that
$R(a^\star)\leq 1/2-\eta$ on admissible held-out bundles, with all leakage and
multiple-comparison corrections applied. No such inequality is presently
reported by this manuscript.
\end{proposition}

\section{Reviewer and Adjudicator Independence Architecture}
\sectionstatus{protocol_defined}

Stage3-2 defined criteria; Stage3-3 now binds those criteria into selection and
assignment packages. This still does not instantiate reviewers or adjudicators,
but it removes the ambiguity about what must exist before such instantiation is
honest.

\section{Reviewer Candidate Pool and Assignment Stub Architecture}
\sectionstatus{protocol_defined}

Stage3-4B deepens the reviewer and adjudicator layer from criteria and workflow
packages into instantiated-empty workflow objects:

\begin{itemize}
\item reviewer candidate pool template;
\item adjudicator candidate pool template;
\item reviewer assignment stub;
\item adjudicator assignment stub;
\item reviewer blinding stub.
\end{itemize}

Each slot remains null-bearing by design. The point is not to simulate people;
the point is to freeze field grammar, slot grammar, null-policy, and the
conditions under which later population becomes honest.

\section{Execution Pilot Package Architecture}
\sectionstatus{protocol_defined}

The execution-pilot package is now a typed family rather than a vague future
bundle. It includes:

\begin{itemize}
\item execution session manifest;
\item collection session schema;
\item reproducibility package skeleton;
\item negative-control bundle skeleton;
\item matched-intervention bundle skeleton;
\item source-target pairing manifest template.
\end{itemize}

\section{Execution Pilot Rehearsal Kit Architecture}
\sectionstatus{protocol_defined}

Stage3-5 adds a rehearsal kit that sits above package readiness and below real
collection. The governing surfaces are:

\begin{itemize}
\item \codestate{docs/execution_pilot_rehearsal_kit.v1.json}
\item \codestate{docs/dry_run_sequence_plan.v1.json}
\item \codestate{docs/rehearsal_event_log_schema.v1.json}
\item \codestate{docs/rehearsal_failure_modes.v1.json}
\item \codestate{docs/rehearsal_reset_policy.v1.json}
\end{itemize}

The rehearsal kit exists to validate sequence integrity, state transitions, and
artifact hygiene without collecting evidence and without emitting decisions.

\section{Reproducibility Skeleton Architecture}
\sectionstatus{protocol_defined}

The reproducibility shell exists as a slot-bearing structure. It records where
later preregistration references, dataset references, control-bundle
references, analysis-shell references, reviewer bundles, and adjudicator
bundles must go.

\section{Dataset Population Boundary}
\sectionstatus{formalized}

Stage3-3 and Stage3-4 enforce a hard distinction between:

\begin{itemize}
\item bundle defined;
\item bundle frozen;
\item bundle populated;
\item bundle adjudicatively interpreted.
\end{itemize}

The current repo has the first state for several families and does not yet have
the latter states.

\section{Dry-Run Event Log Grammar}
\sectionstatus{protocol_defined}

Rehearsal logging is now typed separately from evidence logging. The current
grammar permits:

\begin{itemize}
\item \codestate{rehearsal_initialized};
\item \codestate{sequence_step_checked};
\item \codestate{sequence_step_blocked};
\item \codestate{sequence_step_reset};
\item \codestate{rehearsal_not_started}.
\end{itemize}

None of these event classes may be reinterpreted as human data, reviewer
evidence, or adjudicator action.

\section{Evidence Admissibility Grammar}
\sectionstatus{formalized}

The admissibility grammar still recognizes:

\begin{itemize}
\item independent reviewer report;
\item preregistered human comparator dataset;
\item matched intervention dataset;
\item negative-control dataset;
\item replication log;
\item statistical analysis report;
\item reviewer conflict disclosure;
\item reproducibility package;
\item external failure report;
\item external inconclusive report.
\end{itemize}

Invalid substitutes remain central:

\begin{itemize}
\item internal runtime artifacts alone;
\item self-generated reviewer notes;
\item synthetic human data;
\item LLM-generated ``human'' reports;
\item undocumented manual judgments;
\item unblinded reviewer impressions;
\item post-hoc threshold selection;
\item cherry-picked examples;
\item release inclusion alone;
\item Paper 9 publication alone.
\end{itemize}

\section{Failure, Rejection, and Inconclusive Pathways}
\sectionstatus{protocol_defined}

Three major failure classes remain necessary:

\begin{enumerate}
\item \textbf{protocol invalidation}: preregistration drift, blinding collapse, or conflict violation;
\item \textbf{dataset invalidation}: missing controls, invalid substitution, or insufficient provenance;
\item \textbf{inconclusive adjudication}: admissible evidence exists but remains undecisive under the declared shell.
\end{enumerate}

\section{Statistical Shell Architecture}
\sectionstatus{protocol_defined}

No statistics are executed here. What exists is a declared analysis family
layer, each shell binding channel scope, required inputs, required outputs, and
non-implications.

\section{Adjudication Decision Grammar}
\sectionstatus{protocol_defined}

The future decision family must be typed before it is populated. Paper 10 now
recognizes the following decision tokens as grammar only:

\begin{itemize}
\item \codestate{externally_admissible_but_inconclusive}
\item \codestate{protocol_invalidated}
\item \codestate{adjudication_failed}
\item \codestate{adjudication_passed}
\item \codestate{reviewer_consensus_absent}
\item \codestate{adjudicator_decision_withheld}
\end{itemize}

These tokens are defined as future admissible states. They are not instantiated
anywhere in the current repo as actual outcomes.

\section{State Grammar for Execution Readiness}
\sectionstatus{formalized}

The Stage 3 manuscript surface now distinguishes the following state grammar:

\begin{table}[h]
\centering
\small
\begin{tabular}{L{0.30\textwidth}L{0.58\textwidth}}
\toprule
State & Meaning \\
\midrule
\codestate{schema_defined} & Field grammar exists, but no operational object family is yet instantiated. \\
\codestate{contract_defined} & Allowed relation among objects exists, but no package is assembled. \\
\codestate{package_ready} & The package grammar exists and may be evaluated for completeness. \\
\codestate{instantiated_empty} & The object family exists locally with null-bearing or empty-bearing fields only. \\
\codestate{frozen_local_unfiled} & Local freeze semantics exist, but no external filing is claimed. \\
\codestate{rehearsal_ready} & Rehearsal grammar and rehearsal-kit checks pass, but no live execution shell exists yet. \\
\codestate{live_artifacts_instantiated_unpopulated} & Live operational shells exist, but every collection-facing structure remains empty. \\
\codestate{session_zero_ready} & Session-zero initialization is structurally ready, but still not started. \\
\codestate{collection_openable_under_controls} & Pre-collection admissibility checks pass under contamination and abort controls. \\
\codestate{collection_open_pre_evidential} & Controlled execution may open as a live state without any populated evidence. \\
\codestate{stage3_7_first_collection_openable} & Runtime-health and clean-room audits pass, the opening gate is structurally open, but no collection-start event has been executed. \\
\codestate{collection_started} & A real authorized collection-start event exists. \\
\codestate{evidence_generated} & Real dataset population or admissible log events exist. \\
\codestate{adjudication_result_recorded} & Reviewer-governed or adjudicator-governed decision records exist. \\
\codestate{paper10_final_authorized} & A later separate manuscript-authorization layer has passed. \\
\codestate{dry_run_open} & Dry-run structure may execute without evidence or adjudication. \\
\codestate{externally_filed} & A real external filing reference exists. \\
\codestate{populated_with_admissible_evidence} & The object family is populated by real admissible evidence. \\
\codestate{adjudication_result_bearing} & Reviewer and adjudicator outputs exist and may support result-bearing sections. \\
\bottomrule
\end{tabular}
\end{table}

\section{Controlled Execution Opening}
\sectionstatus{protocol_defined}

Stage3-6 adds a bounded controlled-execution opening layer. This layer is
neither rehearsal-only nor evidence-bearing. It is the first local state in
which live operational artifacts may exist as empty shells while collection
still remains not started.

\begin{itemize}
\item \codestate{execution_opening_policy_defined} means the controlled-execution policy is typed and bound to artifacts.
\item \codestate{collection_openable_under_controls} means pre-collection gates pass under explicit contamination and abort constraints.
\item \codestate{collection_open_pre_evidential} means the system may open a live execution surface while still containing no evidence.
\end{itemize}

\section{Live-But-Unpopulated Artifacts}
\sectionstatus{protocol_defined}

Stage3-6 instantiates a new artifact class: live-but-unpopulated operational
records. These include session-zero manifests, live collection manifests,
empty collection logs, live dataset manifests, live reproducibility manifests,
live decision-record stubs, and live evidence-insertion registers.

Each such surface must explicitly preserve:

\begin{itemize}
\item \codestate{collection_not_started};
\item \codestate{evidence_generation_absent};
\item \codestate{adjudication_not_open};
\item explicit empty markers such as \codestate{empty_dataset = true} where applicable.
\end{itemize}

Therefore live artifact instantiation is not population, and live shells must
never be treated as if they were evidence-bearing records.

\section{Session-Zero Protocol}
\sectionstatus{protocol_defined}

Session-zero is the bounded initialization checkpoint that binds local frozen
record routing, live collection manifests, empty dataset manifests, empty
collection logs, empty decision stubs, contamination controls, abort controls,
and the operator checklist. Session-zero may become ready without collection
starting.

Session-zero remains distinct from:

\begin{itemize}
\item collection start;
\item dataset population;
\item reviewer action;
\item adjudicator action.
\end{itemize}

\section{Pre-Collection Admissibility Gate}
\sectionstatus{protocol_defined}

The pre-collection gate asks whether collection may be opened under controls,
not whether collection has already started. This gate checks session-zero
readiness, emptiness of live artifacts, emptiness of collection logs,
emptiness of decision stubs, absence of evidence, absence of adjudication, and
absence of detected contamination.

Passing this gate yields only \codestate{collection_openable_under_controls}.
It does not yield \codestate{collection_started}.

\section{Contamination, Abort, and Invalidation Controls}
\sectionstatus{protocol_defined}

Stage3-6 adds a more operational invalidation grammar:

\begin{itemize}
\item \textbf{contamination event}: any unauthorized dataset population, evidence-bearing log entry, or decision-bearing record.
\item \textbf{abort event}: any operator-controlled termination of the controlled opening before collection start.
\item \textbf{invalidated session}: any session whose pre-collection boundary, blinding route, or provenance cleanliness is breached.
\end{itemize}

These control classes are meaningful only because Stage3-6 now contains live
artifacts. Without them, contamination and abort would remain abstract policy
tokens rather than execution-opening constraints.

\section{Operator Checklist}
\sectionstatus{protocol_defined}

The operator checklist now matters as a transition object. It binds the local
frozen record route, live manifest route, empty-log route, empty-dataset
state, contamination rules, reset rules, and explicit confirmation that no
evidence exists. Until it is complete, Stage3-6 must not move from
\codestate{session_zero_ready} to \codestate{collection_openable_under_controls}.

\section{Collection-Start Boundary}
\sectionstatus{formalized}

\codestate{collection_open_pre_evidential} is not \codestate{collection_started}.
The boundary is crossed only by a later real operator-authorized collection
start event with provenance, contamination-free status, and a first admissible
log event.

\section{Evidence-Generation Boundary}
\sectionstatus{formalized}

Evidence generation begins only when a real admissible population event enters
the human comparator dataset family, matched intervention family,
negative-control family, collection log, reproducibility package, or
decision-linked provenance chain. No Stage3-6 artifact emitted in this pass
crosses that boundary.

\section{Decision-Record Boundary}
\sectionstatus{formalized}

The live decision-record stub is an execution-opening shell only. It is not a
decision, not an adjudicator action, not a reviewer action, and not an
outcome. The manuscript must preserve that distinction because a populated
decision record belongs to a later evidential and adjudicative state family.

\section{Transition from Session-Zero to First Admissible Collection}
\sectionstatus{blocked_until_execution}

The next real transition is now sharper than in earlier passes:

\begin{enumerate}
\item session-zero becomes ready;
\item the pre-collection gate opens collection under controls;
\item Stage3-6 opens a live but still pre-evidential state;
\item a later explicit collection-start command instantiates the first admissible collection event;
\item only then may the first dataset-population event exist.
\end{enumerate}

\section{Runtime Health Audit}
\sectionstatus{protocol_defined}

Stage3-6B adds a runtime-health audit that asks whether the repository is
actually executable on a bounded local environment rather than only formally
specified. The required audit surfaces include:

\begin{itemize}
\item toolchain visibility;
\item dependency-install behavior under the declared repository install route;
\item repository build behavior;
\item repository test behavior;
\item bounded local-start behavior;
\item verifier-chain behavior for source-of-truth, comparative closure, Stage3-6, Stage3-5, and Paper 9 handoff.
\end{itemize}

The runtime-health audit may pass with warnings. Such warnings may include
engine mismatch, fallback-install compatibility, or dependency advisories. A
warning-bearing pass is still operationally meaningful, but it is not
epistemically elevated into evidence.

\section{Clean-Room Reproducibility}
\sectionstatus{protocol_defined}

Stage3-6B also adds a clean-room reproducibility audit. The clean-room audit
asks whether a bounded local copy can reinstall dependencies, rebuild the
repository, rerun key verifiers, and regenerate Stage 3 artifacts without
relying on hidden machine-specific state. Clean-room reproducibility may pass
with warnings when a compatibility route is required, provided the route is
explicitly documented and the resulting build remains honest.

\section{Dependency and Build Assumptions}
\sectionstatus{formalized}

The runtime-health layer now makes dependency and build assumptions explicit:

\begin{itemize}
\item the declared Node engine remains part of the reproducibility contract;
\item the preferred install route remains \codestate{npm ci};
\item a fallback install route may be recorded only as a compatibility surface;
\item build success must remain distinct from evidence generation;
\item test success must remain distinct from adjudication.
\end{itemize}

\section{Reproducible Artifact Regeneration}
\sectionstatus{protocol_defined}

Stage3-6B adds a source-driven artifact-regeneration audit. This audit checks
whether tracked artifacts can be regenerated from the current source tree and
whether the resulting diffs are:

\begin{itemize}
\item absent;
\item provenance-only;
\item semantically stale but honestly fixed; or
\item blocking.
\end{itemize}

This matters because Stage 3 now depends on a growing artifact family. The
audit therefore separates semantic drift from ordinary lineage and timestamp
drift, and it refuses to treat manual cosmetic rewriting as acceptable
regeneration.

\section{Local Path Leakage Policy}
\sectionstatus{formalized}

Clean-room readiness now requires that live normative surfaces remain free of
machine-specific absolute paths. Historical planning notes or preserved
snapshots may still contain old local references, but they must remain
explicitly classified as historical rather than live normative dependencies.

\section{Controlled Execution Opening Audit}
\sectionstatus{protocol_defined}

The controlled-execution opening layer is now itself audited. Stage3-6B does
not replace Stage3-6; it checks whether Stage3-6 is actually executable,
rebuildable, and reproducible. Therefore the live-but-unpopulated artifacts,
session-zero gate, and pre-collection gate must now survive both runtime
health verification and clean-room reconstruction.

\section{Session-Zero Reproducibility}
\sectionstatus{protocol_defined}

Session-zero reproducibility means that a rebuilt local tree can recreate the
session-zero manifest, operator checklist state, and failure-mode assessment
without any hidden local inputs and without producing collection-start or
evidence-bearing records.

\section{Pre-Collection Gate Reproducibility}
\sectionstatus{protocol_defined}

Pre-collection gate reproducibility means that a rebuilt local tree can
recreate the pre-collection admissibility result, contamination checklist, and
invalidation status while still preserving:

\begin{itemize}
\item \codestate{collection_not_started};
\item \codestate{evidence_generation_absent};
\item \codestate{adjudication_not_open}.
\end{itemize}

\section{Why Reproducibility Does Not Imply Evidence}
\sectionstatus{formalized}

Reproducibility shows that the formal and executable surfaces are rebuildable.
It does not show that any dataset has been populated, that any human evidence
exists, that any reviewer or adjudicator has acted, or that any result-bearing
section is now licensed.

\section{Why Runtime Health Does Not Imply Adjudication}
\sectionstatus{formalized}

Runtime health shows that the code path and verifier path are alive. It does
not imply that collection has started, that evidence has been generated, that
reviewer-governed interpretation exists, or that adjudication is open. A
healthy runtime is still a pre-evidential state until later real execution
events occur.

\section{Transition Conditions for First Admissible Collection}
\sectionstatus{protocol_defined}

Stage3-7 may be classified as \codestate{stage3_7_first_collection_openable}
only when runtime-health and clean-room reproducibility both pass honestly,
path leakage remains clear on live normative surfaces, artifact regeneration is
non-blocking, Stage3-6 remains openable, and Paper 10 remains non-final and
non-result-bearing.

Even then, the current state must still preserve:

\begin{itemize}
\item \codestate{collection_start_not_executed};
\item \codestate{evidence_generation_absent};
\item \codestate{adjudication_not_open};
\item \codestate{decision_record_absent}.
\end{itemize}

\section{Claim Grammar for Eventual Execution-Era Statements}
\sectionstatus{protocol_defined}

Any later admissible result-bearing statement must reference:

\begin{itemize}
\item preregistration identifier;
\item freeze identifier;
\item dataset family identifier;
\item negative-control family identifier;
\item analysis shell identifier;
\item reviewer package identifier;
\item adjudicator package or decision identifier.
\end{itemize}

Forbidden grammar remains explicit:

\begin{itemize}
\item no leap from comparator similarity to human equivalence;
\item no leap from transport-sensitive comparison to identity preservation;
\item no leap from admissibility to metaphysical subjecthood;
\item no leap from dry-run success to public-release recommendation.
\end{itemize}

\section{Theorem, Proposition, and Protocol Mapping}
\sectionstatus{formalized}

\begin{table}[h]
\centering
\small
\begin{tabular}{L{0.22\textwidth}L{0.21\textwidth}L{0.19\textwidth}L{0.18\textwidth}L{0.12\textwidth}}
\toprule
Family & Dependency set & Formalized now & Requires execution & Remains forbidden \\
\midrule
Comparator channel contract & BaseCore, BPF, Stage 2, Stage3-1 & yes & no & result claim \\
Freeze architecture & Stage3-1, Stage3-2, Stage3-3 & yes & no & external filing claim \\
Reviewer package architecture & Stage3-2, Stage3-3 & yes & no & review outcome \\
Pilot package architecture & Stage3-3 & yes & no & evidence claim \\
Dry-run opening grammar & Stage3-3, Stage3-4 & yes & no & adjudication claim \\
Human comparator result section & all above + admissible data & no & yes & human equivalence \\
Adjudication outcome section & all above + reviewer/adjudicator layer + admissible evidence & no & yes & metaphysical closure \\
\bottomrule
\end{tabular}
\end{table}

\section{What Paper 10 Could and Could Not Change}
\sectionstatus{formalized}

Paper 10 is the only layer in the present corpus designed to move evidence out
of the purely internal-support regime. Even a successful execution would not
authorize all possible interpretations.

\begin{table}[h]
\centering
\small
\begin{tabular}{L{0.28\textwidth}L{0.24\textwidth}L{0.34\textwidth}}
\toprule
Claim family & Possible post-execution movement & Still forbidden by default \\
\midrule
Runtime reproducibility & from local/internal to externally reproduced under protocol & consciousness, subjecthood, phenomenality \\
Human comparator similarity & from unpopulated protocol to measured channel comparison & human equivalence or moral parity \\
Bridge predicate discrimination & from BPF scaffold/internal surfaces to adjudicated comparative evidence & phenomenal experience simpliciter \\
Negative-control survival & from internal controls to externally witnessed controls & theorem confirmation beyond stated assumptions \\
Failure or inconclusive result & from open uncertainty to documented non-support or bounded inconclusiveness & rhetorical rescue by post hoc reinterpretation \\
\bottomrule
\end{tabular}
\end{table}

\begin{proposition}[External evidence is not interpretive closure]\label{prop:external-not-closure}
Executed Paper 10 evidence, even if admissible, would not by itself establish
human equivalence, metaphysical subjecthood, biological life, phenomenal
experience, or moral status.
\end{proposition}

\begin{proof}
The Paper 10 protocol measures channel comparison, adjudicator behavior,
negative controls, provenance, contamination, and statistical distinguishability
under frozen rules. None of those objects is identical to the listed
interpretive predicates. Moving a claim from internal support to external
evidence therefore changes its evidential status, not its semantic type.
\end{proof}

\section{Section-Status Ledger}
\sectionstatus{formalized}

\begin{table}[h]
\centering
\small
\begin{tabular}{L{0.34\textwidth}L{0.20\textwidth}L{0.34\textwidth}}
\toprule
Section & Status class & Unlock condition \\
\midrule
Dependency stack & formalized & Update only if upstream stack changes \\
Release-vs-local boundary & formalized & Update only if release authority changes \\
Lineage semantics & formalized & Keep historical-vs-live distinction intact \\
Frozen preregistration architecture & protocol\_defined & Actual local lock still absent \\
Reviewer candidate architecture & protocol\_defined & Selection not started \\
Blinding assignment logic & protocol\_defined & Blinding package not instantiated \\
Execution pilot package architecture & protocol\_defined & Pilot remains unexecuted \\
Reproducibility skeleton architecture & protocol\_defined & Package unpopulated \\
Dataset population boundary & formalized & Keep defined-vs-populated split intact \\
Adjudication decision grammar & protocol\_defined & Decision tokens remain uninstantiated \\
Future result sections & blocked\_until\_execution & Requires admissible execution-era artifacts \\
External adjudication discussion & blocked\_until\_external\_adjudication & Requires admissible reviewer-governed evidence \\
Human equivalence / superiority / metaphysical consequences & forbidden\_claim\_surface & Remains outside current authority surface \\
\bottomrule
\end{tabular}
\end{table}

\section{Evidence Insertion Map}
\sectionstatus{evidence_pending}

\paragraph{Human comparator result insertion point.}
Admissible evidence required: frozen preregistration record, typed human
comparator dataset, negative-control bundle, statistical shell output,
reviewer-selection package, reviewer report family.

\paragraph{External adjudication result insertion point.}
Admissible evidence required: reviewer report set, adjudicator decision record,
blinding log, conflict disclosures, reproducibility package, failure and
inconclusive pathway disclosures.

\paragraph{Transport-sensitive comparison insertion point.}
Admissible evidence required: source-target pairing manifest, transition
lineage map, axis-specific error budgets, transport negative controls,
adjudicator layer.

\paragraph{Instantiated-empty to evidential transition point.}
Admissible evidence required: local frozen record upgraded by real lock event,
real candidate population, real assignment event, real collection session,
real decision record population, and real reproducibility package population.

\section{Open Burdens and Blocked Sections}
\sectionstatus{blocked_until_execution}

\blockedresultstatus

Open burdens remain explicit:

\begin{itemize}
\item real local frozen preregistration lock absent;
\item external filing absent;
\item reviewer selection not started;
\item adjudicator assignment not started;
\item actual blinding assignment absent;
\item admissible human comparator dataset absent;
\item admissible matched-intervention dataset absent;
\item admissible negative-control dataset absent;
\item collection sessions not started;
\item rehearsal event log not started;
\item reproducibility package unpopulated;
\item reviewer and adjudicator decision records absent.
\end{itemize}

\section{Execution-Era Unlock Conditions}
\sectionstatus{protocol_defined}

Before any result-bearing section may become writable, the following must all be
true:

\begin{enumerate}
\item a frozen preregistration record exists;
\item any required external filing state is real and documented;
\item candidate pools are populated by real identities with provenance;
\item reviewer selection is instantiated;
\item adjudicator assignment is instantiated;
\item blinding assignment is instantiated;
\item admissible datasets and control bundles are populated;
\item collection logs and execution sessions are populated by real events;
\item reproducibility package slots are populated;
\item decision grammar is instantiated by real evidence.
\end{enumerate}

\section{Stronger Execution-to-Result Transition Map}
\sectionstatus{blocked_until_execution}

The transition from current Stage3-6 controlled opening to a result-bearing
paper now has six distinct unlock layers:

\begin{enumerate}
\item \textbf{Instantiation layer}: empty records, pools, stubs, and rehearsal kit exist.
\item \textbf{Controlled opening layer}: live operational shells, session-zero, and pre-collection controls exist while collection still remains not started.
\item \textbf{Population layer}: real humans, real sessions, and real evidence populate those empty structures.
\item \textbf{Admissibility layer}: controls, provenance, blinding, and declared shells survive audit.
\item \textbf{Decision layer}: reviewer and adjudicator records become populated by real independent action.
\item \textbf{Manuscript layer}: only then may result-bearing Paper 10 sections be written.
\end{enumerate}

\section{Why Paper 10 Still Cannot Become Final Yet}
\sectionstatus{formalized}

Paper 10 cannot become final because the current repo still lacks the evidence
bearing surfaces that would convert protocol architecture into scientific
claims. The manuscript may now specify the machinery of later freezing,
selection, packaging, dry-run, provenance closure, instantiated-empty records,
and rehearsal-facing governance. It still may not write conclusions whose truth
depends on evidence that does not yet exist.

\section{Future Execution Pathway to Real Stage 3}
\sectionstatus{blocked_until_execution}

The next executable pathway is not ``publish Paper 10.'' It is:

\begin{enumerate}
\item populate the local frozen record honestly;
\item populate reviewer and adjudicator pools honestly;
\item instantiate actual reviewer, adjudicator, and blinding assignments;
\item validate dry-run, rehearsal, and controlled-opening integrity;
\item issue a real collection-start event under Stage3-6 controls;
\item populate admissible datasets, collection logs, and control bundles;
\item run declared statistical shells;
\item bind reviewer and adjudicator outputs;
\item only then write result-bearing sections.
\end{enumerate}

\section{Boundary to Stage 4 Metaphysical Consequence Analysis}
\sectionstatus{forbidden_claim_surface}

Stage 4 remains outside the authority surface of this document. Even later
successful execution-era evidence would still not automatically license
metaphysical subjecthood, moral status, or cross-substrate identity
preservation.

\section{Appendices of Schemas and Templates}
\sectionstatus{protocol_defined}

\subsection*{Appendix A: Frozen Preregistration Record Template}
\sectionstatus{protocol_defined}
Primary reference: \codestate{docs/frozen_preregistration_record_schema.v1.json}.\\
Admissible evidence required before population: actual local lock and, where
required by protocol, real external filing reference.

\subsection*{Appendix B: Human Comparator Preregistration Template}
\sectionstatus{protocol_defined}
Primary reference: \codestate{docs/human_comparator_preregistration_template.v1.json}.\\
Admissible evidence required before population: channel-level freeze and later
dataset lineage.

\subsection*{Appendix C: Reviewer Candidate Schema}
\sectionstatus{protocol_defined}
Primary reference: \codestate{docs/reviewer_candidate_schema.v1.json}.\\
Admissible evidence required before population: real candidate package and
conflict disclosures.

\subsection*{Appendix D: Adjudicator Candidate Schema}
\sectionstatus{protocol_defined}
Primary reference: \codestate{docs/adjudicator_candidate_schema.v1.json}.\\
Admissible evidence required before population: real assignment package and
decision-authority boundary.

\subsection*{Appendix E: Blinding Assignment Template}
\sectionstatus{protocol_defined}
Primary reference: \codestate{docs/blinding_assignment_manifest.v1.json}.\\
Admissible evidence required before population: authorized assignment route and
later unblinding event if one exists.

\subsection*{Appendix F: Execution Pilot Package Manifest}
\sectionstatus{protocol_defined}
Primary reference: \codestate{docs/execution_pilot_package_manifest.v1.json}.\\
Admissible evidence required before population: real session manifests and
later populated bundle references.

\subsection*{Appendix G: Dataset Contract Summary}
\sectionstatus{protocol_defined}
Primary references: \codestate{docs/matched_intervention_dataset_contracts.v1.json}
and \codestate{docs/negative_control_dataset_contracts.v1.json}.\\
Admissible evidence required before result use: populated admissible datasets.

\subsection*{Appendix H: Negative Control Bundle Summary}
\sectionstatus{protocol_defined}
Primary reference: \codestate{docs/pilot_negative_control_bundle.v1.json}.\\
Admissible evidence required before result use: populated bundle and control
admissibility record.

\subsection*{Appendix I: Statistical Shell Summary}
\sectionstatus{protocol_defined}
Primary reference: \codestate{docs/statistical_analysis_shells.v1.json}.\\
Admissible evidence required before result use: shell execution output on real
admissible inputs.

\subsection*{Appendix J: Adjudication Decision Grammar Summary}
\sectionstatus{protocol_defined}
Decision tokens retained as grammar only:
\codestate{externally_admissible_but_inconclusive},
\codestate{protocol_invalidated},
\codestate{adjudication_failed},
\codestate{adjudication_passed},
\codestate{reviewer_consensus_absent},
\codestate{adjudicator_decision_withheld}.

\subsection*{Appendix K: Local Frozen Record Example}
\sectionstatus{protocol_defined}
Primary reference: \codestate{docs/local_frozen_preregistration_record_template.v1.json}.\\
Admissible evidence required before result use: real local lock event and, where
required by route, external filing provenance.

\subsection*{Appendix L: Reviewer Candidate Pool Template}
\sectionstatus{protocol_defined}
Primary reference: \codestate{docs/reviewer_candidate_pool_template.v1.json}.\\
Admissible evidence required before result use: real reviewer identities,
conflict disclosures, and selection events.

\subsection*{Appendix M: Adjudicator Candidate Pool Template}
\sectionstatus{protocol_defined}
Primary reference: \codestate{docs/adjudicator_candidate_pool_template.v1.json}.\\
Admissible evidence required before result use: real adjudicator identity,
assignment provenance, and decision-authority declaration.

\subsection*{Appendix N: Blinding Assignment Stub}
\sectionstatus{protocol_defined}
Primary references: \codestate{docs/blinding_assignment_stub.v1.json} and
\codestate{docs/reviewer_blinding_stub.v1.json}.\\
Admissible evidence required before result use: real assignment route and real
unblinding trace if one later exists.

\subsection*{Appendix O: Execution Session Template}
\sectionstatus{protocol_defined}
Primary reference: \codestate{docs/execution_session_template.v1.json}.\\
Admissible evidence required before result use: real collection session and
real session lineage.

\subsection*{Appendix P: Empty Collection Log Schema}
\sectionstatus{protocol_defined}
Primary reference: \codestate{docs/collection_log_schema.v1.json}.\\
Admissible evidence required before result use: populated log entries from real
collection events.

\subsection*{Appendix Q: Empty Decision Record Template}
\sectionstatus{protocol_defined}
Primary reference: \codestate{docs/decision_record_template.v1.json}.\\
Admissible evidence required before result use: real reviewer and adjudicator
action, not grammar only.

\subsection*{Appendix R: Reproducibility Manifest Example}
\sectionstatus{protocol_defined}
Primary reference: \codestate{docs/reproducibility_package_manifest_empty.v1.json}.\\
Admissible evidence required before result use: populated package slots,
dataset references, shell outputs, and reviewer-facing bundles.

\subsection*{Appendix S: Evidence Insertion Manifest}
\sectionstatus{protocol_defined}
Primary reference: \codestate{docs/evidence_insertion_manifest.v1.json}.\\
Admissible evidence required before result use: actual artifact refs, actual
decision refs, and populated provenance chains.

\subsection*{Appendix T: Rehearsal-Kit Map}
\sectionstatus{protocol_defined}
Primary references: \codestate{docs/execution_pilot_rehearsal_kit.v1.json} and
\codestate{docs/dry_run_sequence_plan.v1.json}.\\
Admissible evidence required before result use: none at rehearsal time, but no
rehearsal artifact may be promoted into evidence.

\subsection*{Appendix U: Dry-Run Event Log Grammar}
\sectionstatus{protocol_defined}
Primary reference: \codestate{docs/rehearsal_event_log_schema.v1.json}.\\
Allowed event classes remain rehearsal-only and do not imply human data,
reviewer action, or adjudicator action.

\subsection*{Appendix V: Controlled Execution Opening Policy}
\sectionstatus{protocol_defined}
Primary reference: \codestate{docs/controlled_execution_opening_policy.v1.json}.\\
Admissible evidence required before result use: none at opening time, but no
controlled-opening artifact may be promoted into evidence by itself.

\subsection*{Appendix W: Session-Zero Protocol and Checklist}
\sectionstatus{protocol_defined}
Primary references: \codestate{docs/session_zero_protocol.v1.json} and
\codestate{docs/session_zero_operator_checklist.v1.json}.\\
Admissible evidence required before result use: a later real collection-start
event and later populated log entries.

\subsection*{Appendix X: Pre-Collection Gate and Contamination Controls}
\sectionstatus{protocol_defined}
Primary references: \codestate{docs/pre_collection_admissibility_gate.v1.json},
\codestate{docs/pre_collection_gate_criteria.v1.json}, and
\codestate{docs/pre_collection_contamination_checklist.v1.json}.\\
Admissible evidence required before result use: a later contamination-free
collection-start event and later admissible dataset population.

\subsection*{Appendix Y: Live-But-Unpopulated Artifact Map}
\sectionstatus{protocol_defined}
Primary references: \codestate{artifacts/external_adjudication/stage3_6/live_collection_manifest.json},
\codestate{artifacts/external_adjudication/stage3_6/live_human_comparator_dataset_manifest.json},
and parallel Stage3-6 live manifests.\\
Allowed interpretation: live operational shell only. Forbidden interpretation:
populated dataset, evidence, or adjudication.

\subsection*{Appendix Z: First Admissible Collection Unlock Map}
\sectionstatus{protocol_defined}
Primary references: \codestate{docs/first_admissible_collection_unlock_map.v1.json}
and \codestate{docs/FIRST_ADMISSIBLE_COLLECTION_UNLOCK_MAP.md}.\\
Admissible evidence required before result use: actual reviewer selection,
actual adjudicator assignment, actual blinding assignment, operator
authorization, collection-start command, and first populated collection-log
event.

\subsection*{Appendix AA: Runtime Health Audit}
\sectionstatus{protocol_defined}
Primary references: \codestate{docs/runtime_health_audit.v1.json},
\codestate{docs/RUNTIME_HEALTH_AUDIT.md}, and
\codestate{artifacts/runtime_health/runtime_health_audit.json}.\\
Allowed interpretation: executable health and verifier-chain status only.
Forbidden interpretation: evidence, collection start, or adjudication.

\subsection*{Appendix AB: Clean-Room Reproducibility Audit}
\sectionstatus{protocol_defined}
Primary references: \codestate{docs/clean_room_reproducibility_audit.v1.json},
\codestate{docs/CLEAN_ROOM_REPRODUCIBILITY_AUDIT.md}, and
\codestate{artifacts/runtime_health/clean_room_reproducibility_audit.json}.\\
Allowed interpretation: bounded rebuildability only. Forbidden interpretation:
external adjudicative support, empirical support, or manuscript finality.

\subsection*{Appendix AC: Local Path Leakage Audit}
\sectionstatus{protocol_defined}
Primary references: \codestate{docs/PATH_AND_LOCAL_LEAKAGE_AUDIT.md} and
\codestate{artifacts/runtime_health/path_and_local_leakage_audit.json}.\\
Allowed interpretation: leak classification only. Forbidden interpretation:
claim that historical snapshots were rewritten or erased.

\subsection*{Appendix AD: Artifact Regeneration Audit}
\sectionstatus{protocol_defined}
Primary references: \codestate{docs/ARTIFACT_REGENERATION_AUDIT.md} and
\codestate{artifacts/runtime_health/artifact_regeneration_audit.json}.\\
Allowed interpretation: regeneration and drift classification only. Forbidden
interpretation: result-bearing evidence insertion.

\subsection*{Appendix AE: Stage3-7 First Collection Opening Criteria}
\sectionstatus{protocol_defined}
Primary references:
\codestate{docs/stage3_7_first_admissible_collection_opening_criteria.v1.json},
\codestate{docs/first_admissible_collection_start_policy.v1.json},
\codestate{artifacts/external_adjudication/stage3_7/stage3_7_opening_gate.json},
and \codestate{artifacts/external_adjudication/stage3_7/first_collection_opening_manifest.json}.\\
Allowed interpretation: first-collection opening criteria only. Forbidden
interpretation: executed collection start, evidence generation, or adjudication.

\section{Blocked Result Surfaces}
\sectionstatus{blocked_until_external_adjudication}

\subsection{Human Comparator Results}
\blockedresultstatus

\paragraph{Admissible evidence required.}
Frozen preregistration record required. Artifact references required. Negative
controls required. Statistical shell required. Reviewer package required.

\subsection{External Adjudication Results}
\blockedresultstatus

\paragraph{Admissible evidence required.}
Independent reviewer reports required. Conflict disclosures required. Blinding
log required. Adjudicator decision record required.

\subsection{Comparative Consequence Claims}
\blockedresultstatus

\paragraph{Admissible evidence required.}
Execution-era artifacts required. Adjudication-facing evidence required.
Identity-preservation, metaphysical, and moral-status claims remain forbidden
even after this section becomes writable.

\section{Non-Claim Ledger}
\sectionstatus{formalized}

This manuscript does not establish consciousness, phenomenality in the strong
sense, human equivalence, superiority, inferiority, metaphysical subjecthood,
cross-substrate identity preservation, or moral status.

\end{document}
